\section{Comparación del modelo con los resultados de prediseño}

\subsection{Período principal de traslación}

El modelo entrega un período fundamental de $T^* = 1{,}259$ s (modo 1), el cual 
corresponde a un modo predominantemente torsional con un 44{,}2\% de participación 
en $R_Z$. Los períodos traslacionales principales son $T_X = 1{,}141$ s en dirección 
X (modo 2) y $T_Y = 0{,}854$ s en dirección Y (modo 3).

El período aproximado calculado en las bases de cálculo fue $T = 1{,}05$ s. La diferencia con el período traslacional 
en dirección X es de aproximadamente un 8.7\%, lo cual es razonable, ya que el resultado del prediseño 
no considera dirección de análisis ni incorpora la rigidez real de los muros 
modelados. 

\subsection{Corte basal}

El corte basal del modelo, tomado en 
el nivel 1, es $Q_{0,X} = 746{,}41$ tonf en dirección X y $Q_{0,Y} = 753{,}54$ tonf 
en dirección Y. El corte basal normativo calculado en las bases de cálculo, considerando 
el coeficiente sísmico mínimo $C_{min} = 0{,}06$, es $Q_0 = 766{,}05$ tonf en 
ambas direcciones. La diferencia es inferior al 3\% en ambas direcciones, lo que 
indica que el modelo es consistente con las bases de cálculo.
